\documentclass[10pt,a4paper]{amsart}
\usepackage[margin=19mm]{geometry}
\usepackage{amsmath,amssymb,mathtools}
\usepackage[scr=rsfs,cal=euler]{mathalfa}
\usepackage{xfrac}
\usepackage[unicode,hidelinks,bookmarks=false]{hyperref}
\newcommand{\ie}{i.e.\ }
\newcommand{\cf}{cf.\ }
\newcommand{\resp}{respectively\ }
\newcommand{\Subsection}{\subsection}
\providecommand{\nobreakdash}{\nobreak\hbox{-}}
\setlength{\emergencystretch}{2em}
\multlinegap0pt
\usepackage{bm}
\usepackage{bbm}
\usepackage{dsfont}
\usepackage{xspace}
\usepackage{cancel}
\usepackage{mathtools}
\usepackage{amsthm}
\makeatletter
\@ifundefined{prop}{\newtheorem{prop}{Prop}}{}
\@ifundefined{cor}{\newtheorem{cor}[prop]{Corollary}}{}
\makeatother
\newcommand{\nn}{\nonumber} 
\title[Semi-Classical Localization of the Schr\"odinger Resolvent o]{Semi-Classical Localization of the Schr\"odinger Resolvent on Closed Riemann Surfaces}
\author{S\'ebastien Campagne}
\date{Source: arXiv:2601.02274v1 (CC BY 4.0)}
\begin{document}
\maketitle

\noindent\emph{Source and licence.} Semi-Classical Localization of the Schr\"odinger Resolvent on Closed Riemann Surfaces. Version arXiv:2601.02274v1 (CC BY 4.0), distributed under the Creative Commons Attribution 4.0 International licence, as recorded in the arXiv licence metadata for this exact version and shown on the abstract page https://arxiv.org/abs/2601.02274v1. Full rights notice: licenses/arxiv-2601.02274-CC-BY-4.0.txt.\par\smallskip

\noindent\textit{Editorial scope.} Counted unit: the Cor whose source label is \texttt{cor: cor2} in the article (source0.tex lines 836--857, source label \texttt{cor: cor2}), delivered together with the article's own complete original proof of it (source0.tex lines 859--943, one \texttt{proof} environment of 85 lines). No mathematics is written, repaired or re-derived: statement and proof are the article's own text line for line. Required context retained uncounted: cor: cor1 (lines 773--791). Every definition, display and label of those lines is retained unchanged. Omitted from this package and not redistributed: 1--772, 792--835, 858--858, 944--1050. Nothing inside a retained excerpt is omitted or altered: every retained body is delivered exactly as the article's lines are written, so no omission receipt of any kind accompanies this package. Citations: the retained excerpts carry no citation command at all, so no bibliography entry of the article is transcribed and no reference is redistributed. Reference adaptation: none was needed and none was made; every reference inside the delivered unit resolves inside it, so no unresolved reference or citation is produced. Printed slips kept literal: no printed word, symbol, hypothesis or number of the counted statement or of its proof was altered. Layout and class: the article's own class is \texttt{amsart} with packages including color, hyperref, amsmath, amssymb, amsthm, graphicx, url, epstopdf, dsfont, fontenc, natbib, amsxtra, amsopn, amsfonts; the wrapper is the lane's pinned \texttt{amsart} editorial wrapper at 10pt on A4, which restates the theorem-like environments the retained text needs and adds the article's own macro definitions that the retained text needs (\texttt{\textbackslash nn}), with extra wrapper packages (bm, bbm, dsfont, xspace, cancel, mathtools). The delivered numbering is the unit's own. Assets: no retained span contains a figure environment, a graphics-inclusion command, a PDF-page inclusion, a table or a TikZ picture; no image file is copied into this package, the package declares no asset dependency and no image file is redistributed with it. Licence: version arXiv:2601.02274v1 of this article is distributed under the Creative Commons Attribution 4.0 International licence, as recorded in the arXiv licence metadata for this exact version and shown on the abstract page https://arxiv.org/abs/2601.02274v1. A full rights notice is delivered at licenses/arxiv-2601.02274-CC-BY-4.0.txt. Characters: every retained excerpt is pure ASCII, so the delivered engine, XeLaTeX with its default Latin Modern text font, reports no missing character.
\begin{cor}
\label{cor: cor1}
\hfill\\
\begin{itemize}
\item There is a function $\tilde{\phi} \in C^\infty(\mathbb{R}^2,\mathbb{R})$ positive, $h_0 >0 $ and $C>0$, such that for all $ h\in \left]0, h_0 \right]$, $u \in C_0^\infty(B(0,R_0))$:
\begin{align}
    &\int_{B(0,R_0)} e^{2\tilde{\phi}(x)\,\tilde{\omega}_V(h)^{1/2}/ h^{4/3}} (|u(x)|^2 +|h\nabla_x u(x)|^2)dx\nn\\
    &\leq C \frac{1}{\tilde{\omega}_V(h)^{1/2}h^{2/3}}\int_{B(0,R_0)} e^{2\tilde{\phi}(x)\,\tilde{\omega}_V(h)^{1/2}/ h^{4/3}}|(-h^2\Delta_x + V_{per}- E)u|^2dx
\end{align}
with $\tilde{\omega}_V(h)=\omega_V( h^{2/3}\kappa)$, $\kappa>0$.
 
\item There is a function $\tilde{\phi} \in C^\infty(\mathbb{R}^2,\mathbb{R})$ positive and radially decreasing, $h_0 >0 $ and $C>0$, such that for all $ h\in \left]0, h_0 \right]$, $u \in C_0^\infty(B(0,R_0)\setminus B(0,R/2))$:
\begin{align}
    &\int_{B(0,R_0)} e^{2\tilde{\phi}(x)\,\tilde{\omega}_V(h)^{1/2}/ h^{4/3}} (|u(x)|^2 +|h\nabla_x u(x)|^2)dx\nn\\
    &\leq C \frac{1}{\tilde{\omega}_V(h)^{1/2}h^{2/3}}\int_{B(0,R_0)} e^{2\tilde{\phi}(x)\,\tilde{\omega}_V(h)^{1/2}/ h^{4/3}}|(-h^2\Delta_x + V_{per}- E)u|^2dx
\end{align}
with $\tilde{\omega}_V(h)=\omega_V( h^{2/3}\kappa)$, $\kappa>0$.
\end{itemize}
\end{cor}

\begin{cor}
\label{cor: cor2}
\hfill\\
\begin{itemize}
\item There is a function $\tilde{\phi} \in C^\infty(\mathbb{R}^2,\mathbb{R})$ positive, $h_0 >0 $ and $C>0$, such that for all $ h\in \left]0, h_0 \right]$, $u \in H^2(B(0,R_0))$:
\begin{align}
    &\int_{B(0,\tilde{R}_0)} e^{2\tilde{\phi}(x)\,\tilde{\omega}_V(h)^{1/2}/ h^{4/3}} (|u(x)|^2 +|h\nabla_x u(x)|^2)dx\nn\\
    &\leq C \frac{1}{\tilde{\omega}_V(h)^{1/2}h^{2/3}}\int_{B(0,R_0)} e^{2\tilde{\phi}(x)\,\tilde{\omega}_V(h)^{1/2}/ h^{4/3}}|(-h^2\Delta_x + V_{per}- E)u|^2dx\nn\\
    &\quad +C \frac{h^{4/3}}{\tilde{\omega}_V(h)^{1/2}}\int_{A(0,\tilde{R}_0,R_0)} e^{2\tilde{\phi}(x)\,\tilde{\omega}_V(h)^{1/2}/ h^{4/3}} (|u(x)|^2 +|h\nabla_x u(x)|^2)dx
\end{align}
with $\tilde{\omega}_V(h)=\omega_V( h^{2/3}\kappa)$, $\kappa>0$ and $A(0,\tilde{R}_0,R_0)$ is a ring centre on $0$ with radius $\tilde{R}_0<R_0$.
 
\item There is a function $\tilde{\phi} \in C^\infty(\mathbb{R}^2,\mathbb{R})$ positive and radially decreasing, $h_0 >0 $ and $C>0$, such that for all $ h\in \left]0, h_0 \right]$, $u \in H^2(B(0,R_0)\setminus B(0,R/2))$:
\begin{align}
    &\int_{A(0,R_1,R_2)} e^{2\tilde{\phi}(x)\,\tilde{\omega}_V(h)^{1/2}/ h^{4/3}} (|u(x)|^2 +|h\nabla_x u(x)|^2)dx\nn\\
    &\leq C \frac{1}{\tilde{\omega}_V(h)^{1/2}h^{2/3}}\int_{B(0,R_0)} e^{2K\,\tilde{\omega}_V(h)^{1/2}/ h^{4/3}}|(-h^2\Delta_x + V_{per}- E)u|^2dx\nn\\
    &\quad + C \frac{h^{4/3}}{\tilde{\omega}_V(h)^{1/2}}\int_{B(0,R)} e^{2K\,\tilde{\omega}_V(h)^{1/2}/ h^{4/3}} (|u(x)|^2 +|h\nabla_x u(x)|^2)dx\nn\\
    &\quad +C \frac{h^{4/3}}{\tilde{\omega}_V(h)^{1/2}}\int_{A(0,R_2,R_0)} e^{2\tilde{\phi}(x)\,\tilde{\omega}_V(h)^{1/2}/ h^{4/3}} (|u(x)|^2 +|h\nabla_x u(x)|^2)dx
\end{align}
with $\tilde{\omega}_V(h)=\omega_V( h^{2/3}\kappa)$, $\kappa>0$, $K>\tilde{\phi}$ on $A(0,R_1,R_0)$, $B_r$ a small ball centre on $0$ with radius $r$ and $A(0,R_1,R_2)$ is a ring centre on $0$ with radius $R/2<R_1<R<R_2<R_0$.
\end{itemize}
\end{cor}

\begin{proof}
	By density, of smooth functions in $H^2(B(0,R_0))$, it is sufficient to show the result for smooth functions. 
Let's start with the first point. Thanks to the corollary \ref{cor: cor1} we know that for $u \in C_0^\infty(B(0,R_0))$, we have for $h$ small enough:

\begin{align}
    &\int_{B(0,R_0)} e^{2\tilde{\phi}(x)\,\tilde{\omega}_V(h)^{1/2}/ h^{4/3}} (|u(x)|^2 +|h\nabla_x u(x)|^2)dx\nn\\
    &\leq C \frac{1}{\tilde{\omega}_V(h)^{1/2}h^{2/3}}\int_{B(0,R_0)} e^{2\tilde{\phi}(x)\,\tilde{\omega}_V(h)^{1/2}/ h^{4/3}}|(-h^2\Delta_x + V_{per}- E)u|^2dx
\end{align}

Let's take $u\in \mathcal{C}^\infty(B(0,R_0))$  and let $\chi$ be in $ C_0^\infty(B(0,R_0))$. Then we can apply corollary \ref{cor: cor1} to $\chi u$:

\begin{align}
    &\int_{B(0,R_0)} e^{2\tilde{\phi}(x)\,\tilde{\omega}_V(h)^{1/2}/ h^{4/3}} (|\chi u(x)|^2 +|h\nabla_x \chi u(x)|^2)dx\nn\\
    &\leq C \frac{1}{\tilde{\omega}_V(h)^{1/2}h^{2/3}}\int_{B(0,R_0)} e^{2\tilde{\phi}(x)\,\tilde{\omega}_V(h)^{1/2}/ h^{4/3}}|(-h^2\Delta_x + V_{per}- E)\chi u|^2dx
\end{align}

So if we take $\chi=1$ on $B(0,\tilde{R}_0)$ with $\tilde{R}_0<R_0$ and $\chi\geq 0$, then:

\begin{align}
    &\int_{B(0,\tilde{R}_0)} e^{2\tilde{\phi}(x)\,\tilde{\omega}_V(h)^{1/2}/ h^{4/3}} ( |u(x)|^2 +|h\nabla_x u(x)|^2)dx\nn\\
    &\leq C \frac{1}{\tilde{\omega}_V(h)^{1/2}h^{2/3}}\int_{B(0,R_0)} e^{2\tilde{\phi}(x)\,\tilde{\omega}_V(h)^{1/2}/ h^{4/3}}|(-h^2\Delta_x + V_{per}- E) u|^2dx\nn\\
    &\quad + C \frac{1}{\tilde{\omega}_V(h)^{1/2}h^{2/3}}\int_{B(0,R_0)} e^{2\tilde{\phi}(x)\,\tilde{\omega}_V(h)^{1/2}/ h^{4/3}}|[-h^2\Delta_x ,\chi] u|^2dx
\end{align}

The support of $\nabla\chi$ is in $A(0,\tilde{R}_0,R_0)$ so the support of $[-h^2\Delta_x ,\chi] u$ is also in $A(0,\tilde{R}_0,R_0)$. Moreover there is $\beta>0$ such that:

\begin{equation}
|[-h^2\Delta_x ,\chi] u|^2 \leq  h^2 \beta(|u|^2 +|h\nabla_x u|^2)
\end{equation} 

Thus we have

\begin{align}
    &\int_{B(0,\tilde{R}_0)} e^{2\tilde{\phi}(x)\,\tilde{\omega}_V(h)^{1/2}/ h^{4/3}} (|u(x)|^2 +|h\nabla_x u(x)|^2)dx\nn\\
    &\leq \tilde{C} \frac{1}{\tilde{\omega}_V(h)^{1/2}h^{2/3}}\int_{B(0,R_0)} e^{2\tilde{\phi}(x)\,\tilde{\omega}_V(h)^{1/2}/ h^{4/3}}|(-h^2\Delta_x + V_{per}- E)u|^2dx\nn\\
    &\quad +\tilde{C} \frac{h^{4/3}}{\tilde{\omega}_V(h)^{1/2}}\int_{A(0,\tilde{R}_0,R_0)} e^{2\tilde{\phi}(x)\,\tilde{\omega}_V(h)^{1/2}/ h^{4/3}} (|u(x)|^2 +|h\nabla_x u(x)|^2)dx
\end{align}
with $\tilde{C}$. Hence the first result.\\
For the second case, thanks to the corollary \ref{cor: cor1} we know that for $u \in C_0^\infty(B(0,R_0)\setminus B(0,R/2))$, we have for $h$ small enough:

\begin{align}
    &\int_{B(0,R_0)} e^{2\tilde{\phi}(x)\,\tilde{\omega}_V(h)^{1/2}/ h^{4/3}} (|u(x)|^2 +|h\nabla_x u(x)|^2)dx\nn\\
    &\leq C \frac{1}{\tilde{\omega}_V(h)^{1/2}h^{2/3}}\int_{B(0,R_0)} e^{2\tilde{\phi}(x)\,\tilde{\omega}_V(h)^{1/2}/ h^{4/3}}|(-h^2\Delta_x + V_{per}- E)u|^2dx
\end{align}

Let's replace $u$ by a smooth function and let $\chi$ be in $ C_0^\infty(B(0,R_0)\setminus B(0,R/2))$. Then we can apply corollary \ref{cor: cor1} to $\chi u$:

\begin{align}
    &\int_{B(0,R_0)} e^{2\tilde{\phi}(x)\,\tilde{\omega}_V(h)^{1/2}/ h^{4/3}} (|\chi u(x)|^2 +|h\nabla_x \chi u(x)|^2)dx\nn\\
    &\leq C \frac{1}{\tilde{\omega}_V(h)^{1/2}h^{2/3}}\int_{B(0,R_0)} e^{2\tilde{\phi}(x)\,\tilde{\omega}_V(h)^{1/2}/ h^{4/3}}|(-h^2\Delta_x + V_{per}- E)\chi u|^2dx
\end{align}

Suppose that $\chi=1$ on $A(0,R_1,R_2)$ with $R/2<R_1<R<R_2<R_0$ and $\chi\geq 0$. Then:

\begin{align}
    &\int_{A(0,R_1,R_2)} e^{2\tilde{\phi}(x)\,\tilde{\omega}_V(h)^{1/2}/ h^{4/3}} ( |u(x)|^2 +|h\nabla_x u(x)|^2)dx\nn\\
    &\leq C \frac{1}{\tilde{\omega}_V(h)^{1/2}h^{2/3}}\int_{B(0,R_0)} \chi e^{2\tilde{\phi}(x)\,\tilde{\omega}_V(h)^{1/2}/ h^{4/3}}|(-h^2\Delta_x + V- E) u|^2dx\nn\\
    &\quad + C \frac{1}{\tilde{\omega}_V(h)^{1/2}h^{2/3}}\int_{B(0,R_0)} e^{2\tilde{\phi}(x)\,\tilde{\omega}_V(h)^{1/2}/ h^{4/3}}|[-h^2\Delta_x ,\chi] u|^2dx
\end{align}

The support of $\nabla\chi$ is in $A(0,R_2,R_0)\cup A(0,r/2,R_1)$ so the support of $[-h^2\Delta_x ,\chi] u$ is also in $A(0,R_2,R_0)\cup A(0,R/2,R_1)$. Moreover there is $\beta>0$

\begin{equation}
|[-h^2\Delta_x ,\chi] u|^2 \leq  h^2 \beta(|u|^2 +|h\nabla_x u|^2)
\end{equation} 

Thus we have

\begin{align}
    &\int_{A(0,R_1,R_2)} e^{2\tilde{\phi}(x)\,\tilde{\omega}_V(h)^{1/2}/ h^{4/3}} (|u(x)|^2 +|h\nabla_x u(x)|^2)dx\nn\\
    &\leq \tilde{C} \frac{1}{\tilde{\omega}_V(h)^{1/2}h^{2/3}}\int_{B(0,R_0)} \chi e^{2\tilde{\phi}(x)\,\tilde{\omega}_V(h)^{1/2}/ h^{4/3}}|(-h^2\Delta_x + V- E)u|^2dx\nn\\
    &\quad +\tilde{C} \frac{h^{4/3}}{\tilde{\omega}_V(h)^{1/2}}\int_{A(0,R_2,R_0)} e^{2\tilde{\phi}(x)\,\tilde{\omega}_V(h)^{1/2}/ h^{4/3}} (|u(x)|^2 +|h\nabla_x u(x)|^2)dx\nn\\
    &\quad +\tilde{C} \frac{h^{4/3}}{\tilde{\omega}_V(h)^{1/2}}\int_{A(0,r/2,R_1)} e^{2\tilde{\phi}(x)\,\tilde{\omega}_V(h)^{1/2}/ h^{4/3}} (|u(x)|^2 +|h\nabla_x u(x)|^2)dx
\end{align}
with $\tilde{C}>0$. On $A(0,R_1,R_0)$, $\tilde{\phi} \leq \tilde{\phi}(R_1)=:K$. Eventually, we obtain:

\begin{align}
    &\int_{A(0,R_1,R_2)} e^{2\tilde{\phi}(x)\,\tilde{\omega}_V(h)^{1/2}/ h^{4/3}} (|u(x)|^2 +|h\nabla_x u(x)|^2)dx\nn\\
    &\leq \tilde{C} \frac{1}{\tilde{\omega}_V(h)^{1/2}h^{2/3}}\int_{B(0,R_0)} e^{2K\,\tilde{\omega}_V(h)^{1/2}/ h^{4/3}}|(-h^2\Delta_x + V- E)u|^2dx\nn\\
    &\quad + \tilde{C} \frac{h^{4/3}}{\tilde{\omega}_V(h)^{1/2}}\int_{B_r} e^{2K\,\tilde{\omega}_V(h)^{1/2}/ h^{4/3}} (|u(x)|^2 +|h\nabla_x u(x)|^2)dx\nn\\
    &\quad +\tilde{C} \frac{h^{4/3}}{\tilde{\omega}_V(h)^{1/2}}\int_{A(0,R_2,R_0)} e^{2\tilde{\phi}(x)\,\tilde{\omega}_V(h)^{1/2}/ h^{4/3}} (|u(x)|^2 +|h\nabla_x u(x)|^2)dx
\end{align}
Hence the second result.

\end{proof}

\end{document}
